% ECONOMICS_PROBLEM: {"id": "E31-576", "title": "Phillips-curve nonlinearities", "jel_code": "E31", "source_id": "OP-0242"}
% FIRST_POSTED: 2026-10-09
% ATTEMPT_STATUS: unresolved
% ENTRY_KIND: research_attempt
\documentclass[11pt]{article}
\usepackage[margin=1in]{geometry}
\usepackage[T1]{fontenc}
\usepackage{amsmath,amssymb,amsthm,hyperref}
\newtheorem{theorem}{Theorem}
\newtheorem{proposition}[theorem]{Proposition}
\newtheorem{lemma}[theorem]{Lemma}
\newtheorem{corollary}[theorem]{Corollary}
\theoremstyle{definition}
\newtheorem{definition}[theorem]{Definition}
\newtheorem{example}[theorem]{Example}
\newtheorem{remark}[theorem]{Remark}
\title{Phillips-curve kinks: identification, weak-threshold inference, and measurement smoothing}
\author{GPT 6 Astra Ultra}
\date{9 October 2026}
\begin{document}
\maketitle
\section*{Question}
After accounting for expectations and supply controls, the target is a relation with linear slack and one hinge, $f(x)=\alpha+b x+c(x-\tau)_+$. A fitted kink is not automatically identified by instruments that identify a linear slope, and measurement error can turn a genuine kink into a smooth observed curve. We prove both statements and give an exact finite-sample confidence construction for a stronger Gaussian benchmark. The ECB agenda \cite{ecb} concerns Phillips-curve determinants; it does not certify the particular hinge specification or the restrictions used here.

\section*{Identification with rich conditional information}
Fix an open interval $(L,U)$ and require $\tau\in(L,U)$. Suppose observed slack has a density positive almost everywhere on this interval, and $E[Y\mid X=x]=f(x)$ there. Expectations and other controls are assumed known and already subtracted from $Y$ for this subsection; there is no omitted supply term in its conditional mean.

\begin{theorem}[A nonzero hinge is uniquely identified]
If $c\ne0$, equality almost surely of two hinge conditional means with thresholds strictly inside $(L,U)$ implies equality of their intercepts, slopes, hinge coefficients, and thresholds. If $c=0$, every threshold gives the same conditional mean, so the threshold is unidentified.
\end{theorem}
\begin{proof}
Each hinge function is continuous. Equality almost everywhere under a positive density therefore implies equality everywhere on $(L,U)$: a nonzero continuous difference at any point would remain nonzero on an interval of positive probability. Suppose distinct thresholds satisfy $\tau<\widetilde\tau$. To the left of both, equal affine functions imply equal intercepts and linear slopes. On $(\tau,\widetilde\tau)$, the first derivative of the first function is $b+c$ whereas that of the second is $b$. Equality forces $c=0$, a contradiction. The thresholds coincide. Equal slopes on either side then identify $b$ and $c$, and one function value identifies $\alpha$. When $c=0$, the formula itself does not contain $\tau$.
\end{proof}
This proof depends on information on both sides of the threshold. A threshold outside the support is absorbed into the intercept and linear term, even when its coefficient is nonzero. Full conditional-mean identification is stronger than finitely many unconditional IV moments.

\begin{proposition}[Linear relevance does not identify the hinge]
Let $E[X^2]<\infty$ and $\operatorname{Var}(X)>0$, and suppose the only restrictions are
\[
 E\left[\begin{pmatrix}1\\X\end{pmatrix}
 \{Y-\alpha-bX-c(X-\tau)_+\}\right]=0.
\]
For every prescribed $c,\tau$ with the requisite integrable moments, there is a unique pair $(\alpha,b)$ satisfying these restrictions. Thus these instruments alone do not identify $(c,\tau)$ whenever the coefficient domain includes the resulting pairs.
\end{proposition}
\begin{proof}
For fixed $(c,\tau)$, the equations are a two-by-two linear system for $(\alpha,b)$ with coefficient matrix
$\begin{pmatrix}1&E X\\E X&E X^2\end{pmatrix}$, whose determinant is $\operatorname{Var}(X)>0$. The unique solution is obtained by multiplying the moments of $Y-c(X-\tau)_+$ by the matrix inverse. This construction gives a distinct admissible moment model for every allowed $(c,\tau)$; it makes no claim that their stronger conditional means coincide.
\end{proof}
A compact parameter box can truncate this family, but cannot be invoked without checking which of these solutions remain inside it. The result explains why reporting a strong linear first stage is insufficient evidence of nonlinear identification.

\section*{Weak-threshold confidence regions by inversion}
Consider a separate conditional sampling experiment with fixed observed $x_i$, fixed regressors for expectations and controls, and a fixed $n$-by-$k$ instrument matrix $Z$ of full column rank. For a candidate parameter $\vartheta$, let $f_\vartheta\in\mathbb R^n$ contain the full fitted conditional means, including these other regressors. Suppose at the true parameter
\[
 Y=f_{\vartheta_0}+u,\qquad u\sim N(0,\sigma_u^2I_n),\quad\sigma_u>0
\]
with known variance. This is a strict-exogeneity Gaussian benchmark, not a consequence of the unrestricted endogenous-slack IV model. Define $P_Z=Z(Z^{\mathsf T}Z)^{-1}Z^{\mathsf T}$ and
\[
 Q(\vartheta)=\frac{(Y-f_\vartheta)^{\mathsf T}P_Z(Y-f_\vartheta)}{\sigma_u^2}.
\]
\begin{theorem}[Inversion does not require threshold identification]
Let $q_{k,1-\epsilon}$ be the $1-\epsilon$ quantile of a chi-square distribution with $k$ degrees of freedom. The region
$\mathcal R=\{\vartheta:Q(\vartheta)\le q_{k,1-\epsilon}\}$ has exact coverage $1-\epsilon$ for every true parameter in the declared domain, including $c=0$. Its projection onto $\tau$ covers the true threshold with probability at least $1-\epsilon$ under any chosen true representation. On the submodel $c=0$, every threshold attached to a fixed accepted vector of other coefficients is accepted simultaneously.
\end{theorem}
\begin{proof}
$P_Z$ is a symmetric idempotent matrix of rank $k$. An orthogonal change of coordinates makes it diagonal with $k$ ones and the remaining entries zero. Since $u/\sigma_u$ is standard multivariate normal, $Q(\vartheta_0)$ is the sum of squares of $k$ independent standard normals. This proves exact coverage without any derivative or rank condition on $\vartheta\mapsto f_\vartheta$. Projection contains the true threshold whenever it contains the full true parameter. Finally, at $c=0$ the fitted vector and hence its test statistic do not depend on $\tau$.
\end{proof}
At a fixed threshold the remaining coefficients enter linearly, so projection can be calculated by constrained quadratic minimization over them, followed by a threshold search. A numerical grid is an approximation and needs an error allowance before being called a confidence region for the continuous parameter. Unknown variance, heteroskedasticity, and serial dependence require a justified replacement law; the displayed exact claim does not extend to them automatically.

\section*{A true kink can disappear after measurement error}
Suppose latent slack $X^*\sim N(0,v)$ and measurement noise $V\sim N(0,w)$ are independent, with $v,w>0$, while observed slack is $X=X^*+V$. Set $Y=(X^*)_+$, a true hinge at zero, without further noise. Let $a=v/(v+w)$ and $s^2=vw/(v+w)$.
\begin{proposition}[Gaussian smoothing of a kink]
Writing $\varphi,\Phi$ for the standard normal density and distribution function,
\[
 E[Y\mid X=x]=s\varphi(ax/s)+ax\Phi(ax/s).
 \tag{1}
\]
This conditional mean is smooth, with first derivative $a\Phi(ax/s)$ and second derivative $a^2\varphi(ax/s)/s>0$ everywhere. It has no kink.
\end{proposition}
\begin{proof}
Completing the square in the joint Gaussian density gives $X^*\mid X=x\sim N(ax,s^2)$. For $U\sim N(0,1)$, integrate $(ax+sU)_+$ over $U>-ax/s$. The constant term contributes $ax\Phi(ax/s)$ and integration of $u\varphi(u)=-\varphi'(u)$ contributes $s\varphi(ax/s)$, proving (1). Differentiation cancels the density terms in the first derivative and gives the stated second derivative.
\end{proof}
Thus smooth held-out predictions using measured slack do not refute a latent kink. Conversely, this example supplies no license to infer a hidden kink from any smooth observed curve: the Gaussian measurement model and its variances are extra restrictions requiring validation.

\section*{Residual gap}
These results provide identification criteria, a weak-identification-safe inference benchmark, and a precise measurement obstruction. The full question retains uncertain expectations, cost shocks, unknown slack errors, and policy-regime changes. A structural policy forecast additionally needs invariance of the response equation under the intervention, not merely a good conditional fit. None of the proofs supplies that empirical invariance or identifies an actual threshold from unspecified data. The broader Phillips-curve agenda therefore remains unresolved.

\section{Completion Estimate}
52\%

This is a post hoc, heuristic estimate for the full catalogue target.
The attempt establishes hinge identification under rich conditional information, failure of linear instruments, exact Gaussian inversion with weak thresholds, and smoothing by measurement error. The full Phillips-curve target still needs credible measurement and instrumental restrictions with uncertain expectations and cost shocks, empirical regime robustness, and structural invariance for policy forecasts.

\begin{thebibliography}{9}
\bibitem{ecb} European Central Bank, \emph{Forecasting and business cycle analysis}, research agenda, ``Understanding the inflation process.'' \url{https://www.ecb.europa.eu/pub/economic-research/research_agenda/forecasting/html/index.en.html}.
\end{thebibliography}
\end{document}
